\chapter{IoT Projects}

\section{Some IoT Projects Examples}

\begin{multicols}{2}
\textbf{Smart Poultry and Fish Farming System}
\begin{enumerate}
   \item NodeMCU (ESP8266)
   \item LCD Display
   \item DHT11/DHT22
   \item Relay Module
   \item SMPS
   \item Ultrasonic Sensor
   \item Moisture Sensor
   \item pH Sensor
   \item TDS Sensor
   \item Turbidity Sensor
   \item MQ-135 Gas Sensor
   \item Water Pump
   \item Light Dependent Resistor (LDR)
   \item Servo Motor
\end{enumerate}

% \begin{table}[H]
% % \renewcommand{\arraystretch}{1.3}
% \begin{tabular}{ll}
%    \textbf{Component} & \textbf{Usage in the project} \\

%    {NodeMCU - ESP8266} & : Microcontroller for processing sensor data and controlling actuators. \\

%    {LCD Display} & : To display sensor readings/data. \\

%    {DHT11/DHT22} & : To monitor temperature and humidity. \\

%    {Relay Module} & : To control on/off switching. \\

%    {SMPS} & : To convert AC 220V to DC 5V for power supply. \\

%    {Ultrasonic Sensor} & : To monitor food level. \\

%    {Moisture Sensor} & : To monitor the moisture of the water level. \\

%    {pH Sensor} & : To monitor water pH (ideal pH: 6.5-8.5) to maintain fish health. \\

%    {TDS Sensor} & : To monitor water quality by measuring some factors. \\

%    {Turbidity Sensor} & : To monitor water clarity in fish farming, as high turbidity indicates dirty water. \\

%    {MQ-135 Gas Sensor} & : To detect harmful gases like ammonia and CO2 to prevent diseases. \\

%    {Water Pump} & : To supply water automatically when required. \\

%    {Light Dependent Resistor (LDR)} & : To monitor light intensity and maintaining a proper lighting. \\

%    {Servo Motor} & : To control automatic feeder for precise feeding schedules. \\

%    {Cooling Fan / Exhaust Fan} & : To control temperature or humidity automatically if required. \\

%    {Mobile/Desktop App} & : To enable real-time remote monitoring and controlling. \\

% \end{tabular}
% \end{table} 

\textbf{Obstacle Avoider Robot}
\begin{enumerate}
   \item Arduino UNO / ESP8266
   \item Motor Driver
   \item Servo Motor
   \item Ultrasonic sensor
   \item Gear Motor
   \item 18650 Battery
\end{enumerate}

\textbf{Line Follower Robot}
\begin{enumerate}
   \item Arduino UNO / ESP8266
   \item Gear Motor
   \item Motor Driver
   \item IR sensor
   \item 18650 Battery
\end{enumerate}

\columnbreak
\textbf{Smart Plant Monitoring System}
\begin{enumerate}
   \item Arduino UNO / ESP8266
   \item PIR Motion sensor
   \item DHT11 Temperature sensor
   \item Soil Moisture sensor
   \item Relay module
   \item 18650 Battery
\end{enumerate}

\textbf{Metal Detector Robot}
\begin{enumerate}
   \item Arduino UNO / ESP8266
   \item Gear Motor
   \item Motor Driver
   \item Metal Detector
   \item LCD display
   \item Joystick Module
   \item 18650 Battery
\end{enumerate}

\textbf{Voice Controlled Car}
\begin{enumerate}
   \item Arduino UNO / ESP8266
   \item Gear Motor
   \item Servo Motor
   \item Motor Driver
   \item Ultrasonic sensor
   \item Bluetooth sensor
   \item 18650 Battery
\end{enumerate}
\end{multicols}

\pagebreak
\subsection{Smart Agriculture}

\textbf{\textit{Components}:}\\[0.5em]
\begin{tabular}{lll}
   1. ESP8266
   &:& Micro-controller for IoT connectivity. \\
   2. PIR Motion sensor
   &:& To detect motion in the farm. \\
   3. DHT11 Temperature sensor
   &:& To monitor temperature and humidity. \\
   4. Soil Moisture sensor
   &:& To measure soil moisture levels. \\
   5. Relay module
   &:& To control devices like pumps or lights. \\
   6. 18650 Battery
   &:& Power supply for the system. \\
   7. LDR (Light Dependent Resistor)
   &:& To monitor light intensity for plant growth. \\
   8. Water Pump
   &:& To irrigate plants based on soil moisture levels. \\
   9. pH Sensor
   &:& To monitor soil pH for optimal plant growth. \\
   10. OLED Display
   &:& To display sensor readings and system status locally. \\
\end{tabular}

\vspace{2em}
\textbf{\textit{Pseudo Code}:}
\begin{pseudocode}   
BEGIN
   INITIALIZE the sensors and actuators

   LOOP FOREVER:
      === Sensor Readings ===
      READ temperature, humidity from DHT11
      READ soilMoistureValue from Soil Moisture Sensor
      READ pHValue from pH Sensor
      READ lightLevel from LDR
      READ motionDetected from PIR Sensor (HIGH/LOW)

      === Irrigation Logic ===
      IF soilMoistureValue < THRESHOLD:
         TURN ON Relay (activate water pump)
         WAIT for irrigationTime seconds
         TURN OFF Relay
      ELSE:
         KEEP water pump OFF

      === Motion Detection Logic ===
      IF motionDetected is HIGH:
         LOG or DISPLAY "Motion detected"
         Send alert to user

      === Display Logic ===
      CLEAR OLED screen
      DISPLAY(temperature,humidity,soil_moisture,pH)

      DELAY 1000
   END LOOP
END
\end{pseudocode}

\pagebreak
\textbf{\textit{Iteration}:}
\begin{arduino}
void loop() {
   // Sensor Readings
   float temperature = dht.readTemperature();
   float humidity = dht.readHumidity();
   int soilMoistureValue = analogRead(soilMoisturePin);
   float pHValue = readPH();
   int lightLevel = analogRead(ldrPin);
   bool motionDetected = digitalRead(pirPin) == HIGH ? true : false;
   
   // Irrigation Logic
   if (soilMoistureValue < moistureThreshold) {
      digitalWrite(WATER_PUMP_RELAY, HIGH);
      delay(irrigationTime * 1000);
      digitalWrite(WATER_PUMP_RELAY, LOW);
   } else {
      digitalWrite(WATER_PUMP_RELAY, LOW);
   }
   
   // Motion Detection Logic
   if (motionDetected) {
      Serial.println("Motion detected!");
   }
   
   // Display Logic
   oled.clearDisplay();
   oled.setCursor(0, 0);
   oled.print("Temp: "); oled.println(temperature);
   oled.print("Humidity: "); oled.println(humidity);
   oled.print("Soil: "); oled.println(soilMoistureValue);
   oled.print("pH: "); oled.println(pHValue);
   oled.print("Light: "); oled.println(lightLevel);
   oled.print("Motion: "); oled.println(motionDetected ? "Yes" : "No");
   oled.display();
   
   delay(1000);
}
\end{arduino}

\pagebreak
\subsection{2WD and 4WD Robot}
\begin{minipage}[t]{0.6\textwidth}
\textbf{\textit{Pseudo Code}:}
\begin{pseudocode}
LOOP FOREVER:
   READ leftSensor and rightSensor

   IF leftSensor IS WHITE AND rightSensor IS WHITE THEN
      MOVE forward
   ELSE IF leftSensor IS BLACK AND rightSensor IS WHITE THEN
      TURN right
   ELSE IF leftSensor IS WHITE AND rightSensor IS BLACK THEN
      TURN left
   ELSE
      STOP motors
END LOOP
\end{pseudocode}
\end{minipage}\hfill
\begin{minipage}[t]{0.3\textwidth}
   Black lines reflect less IR.\\
   So, IR sensors return LOW when over black and HIGH when over white.
   
   \vspace{2em}
   HIGH + LOW = Forward\\
   LOW + HIGH = Backward\\
\end{minipage}


\textbf{\textit{Iteration}:}
\begin{arduino}
void loop() {
   int RS = digitalRead(IR_RIGHT_PIN);
   int LS = digitalRead(IR_LEFT_PIN);

   // Both sensors see white => Move forward
   if (LS == 1 && RS == 1) {
      setMotors(HIGH, HIGH);
   }

   // Right sensor sees black, turn left
   if (LS == 1 && RS == 0) {
      setMotors(LOW, HIGH);
   }

   // Left sensor sees black, turn right
   if (LS == 0 && RS == 1) {
      setMotors(HIGH, LOW);
   }

   // Both sensors see black, stop
   if (LS == 0 && RS == 0) {
      stopMotors();
   }
}

// Combined control function
const setMotors(bool isLeftForward, bool isRightForward) {
   digitalWrite(in1, isLeftForward); digitalWrite(in2, !isLeftForward);
   digitalWrite(in3, isRightForward); digitalWrite(in4, !isRightForward);
   analogWrite(enA, MAX_SPEED); analogWrite(enB, MAX_SPEED);
}

const stopMotors() {
   analogWrite(enA, 0);
   analogWrite(enB, 0);
}
\end{arduino}